\documentclass[12pt]{article}

\usepackage[utf8]{inputenc}
\usepackage[greek,english]{babel}
\usepackage{alphabeta}
\usepackage{amsmath,amssymb}
\usepackage{geometry}
\geometry{a4paper, margin=1in}

\begin{document}

\begin{center}
{\Large \textbf{Γραπτές προαγωγικές εξετάσεις περιόδου Μαΐου – Ιουνίου 2026}}

\vspace{0.5cm}

στο μάθημα της \underline{\hspace{5cm}}

\end{center}

\vspace{1cm}

\textbf{ΗΜΕΡΟΜΗΝΙΑ:} \underline{\hspace{3cm}}

\vspace{0.3cm}

\textbf{ΟΝΟΜΑΤΕΠΩΝΥΜΟ:} \underline{\hspace{8cm}}

\vspace{0.3cm}

\textbf{ΤΑΞΗ:} \underline{\hspace{2cm}} Λυκείου

\vspace{0.3cm}

\textbf{ΕΙΣΗΓΗΤΕΣ:}

\vspace{0.3cm}

\textbf{ΕΠΙΤΗΡΗΤΗΣ:} \underline{\hspace{6cm}}

\vspace{1cm}

\section*{ΘΕΜΑ Α}

\textbf{Α1.} Να γράψετε το τετράδιό σας το νούμερο της πρότασης και δίπλα τη λέξη «Σωστό» αν η πρόταση είναι σωστή ή τη λέξη «Λάθος» αν η πρόταση είναι λανθασμένη.

\hfill \textbf{(Μονάδες __)}

\vspace{0.5cm}

\textbf{Α2.} Να αποδείξετε ότι:

\hfill \textbf{(Μονάδες __)}

\vspace{1cm}

\section*{ΘΕΜΑ Β}

\textbf{Β1.}

\vspace{1cm}

\textbf{Β2.}

\vspace{1cm}

\textbf{Β3.}

\newpage

\section*{ΘΕΜΑ Γ}

\textbf{Γ1.}

\vspace{1cm}

\textbf{Γ2.}

\vspace{1cm}

\textbf{Γ3.}

\vspace{1cm}

\section*{ΘΕΜΑ Δ}

\textbf{Δ1.}

\vspace{1cm}

\textbf{Δ2.}

\vspace{1cm}

\textbf{Δ3.}

\end{document}